\documentclass[a4paper]{article}
\usepackage[polish]{babel}
\usepackage[cp1250]{inputenc}
\usepackage[T1]{fontenc}
\usepackage[dvips]{graphicx}

\pagestyle{headings}
\textwidth      15.5cm
\oddsidemargin    .1cm
\evensidemargin   .1cm

\begin{document}

\thispagestyle{empty}

\begin{minipage}{5cm}
\includegraphics[scale=0.7]{salomon.eps}
\end{minipage}
\begin{minipage}{10cm}
\begin{center}
{\Large Salomon\\
System przetwarzania wiedzy\\}
\end{center}
\end{minipage}

\vspace*{0.5cm}

\hrulefill

\vspace*{1cm}


\begin{flushright}
{\Large \emph{Iteration Description --- C3, ver 1.1} }
\end{flushright}



\begin{flushleft}
{ \Large Historia wersji }\[\]
\begin{tabular}{|p{2.5cm}||p{1.5cm}|p{6cm}|p{3.5cm}|}
  \hline
  Data & Wersja & Opis & Autor \\
  \hline \hline
  27.01.2006 & 1.0 & utworzenie & \emph{Leszek Grzanka}\\
  \hline
  24.02.2006 & 1.1 & drobne korekty & \emph{Nikodem Jura}\\
\end{tabular}
\end{flushleft}

\newpage


\section{Zadania dla iteracji } 

Budowa kompletnej dokumentacji systemu, w szczeg�lno�ci: dokumenty zgodne z metodologi� RUP, dokument�w przydatnych u�ytkownikowi (podr�cznik), oraz dokument�w dost�pnych w sieci wraz z ostateczn� wersj� aplikacji.


\section{Produkt iteracji} 
Dokumentacja projektu, dost�pna w wersji elektronicznej na stronie internetowej produktu.

\section{Plan test�w} 
Sprawdzenie zgodno�ci dokumentacji z wymaganiami stawianymi jej przez metodologi�.

\section{Efekt wykonania iteracji}

Efektem by�o:
\begin{itemize}
 	\item przygotowanie odpowiednich dokument�w wed�ug metodologii RUP (Vision, SDP, SAD, ID, \ldots)
	\item podr�cznik u�ytkownika	
	\item strona WWW produktu, zawieraj�ca pe�n� dokumentacj� oraz finaln� wersj� produktu
	\item system G-Forge umo�liwiaj�cy nadzorowanie procesu powstawania aplikacji
\end{itemize}


\end{document}

